\summarychapter{43}{Rootless Immigrants in a Sick Society}{American Evangelicals}

\begin{inbrief}
In the 1970s conservative Protestantism in the United States turned out to be large, confident and
growing, and still it counted for little in politics, the universities and the press. To explain the
gap the chapter looks back to the decades after the Civil War, when revivalist Protestants gave up
social reform, to the fundamentalist controversy and to the humiliation of William Jennings Bryan. It
then describes the recovery after 1945 around Billy Graham and the spread of Pentecostal experience
into the older churches.
\question{Evangelicals could give millions a personal faith. Could they also renew the public ideals
of the nation, as earlier awakenings had done?}
\end{inbrief}

The chapter begins with the election of 1976, in which the born-again faith of Jimmy Carter became a
campaign topic. Watergate was recent, and many voters wanted moral leadership above all. Others
thought religion had no business in an election. Arthur M.~Schlesinger Jr., the historian who had
advised Kennedy and Johnson, dismissed Carter's talk of closeness to God: ``If he feels that, bully
for him, but it's totally irrelevant to the campaign!'' As a statement about qualifications he was
correct, because the Constitution excludes religious tests. Yet convinced Christians cannot regard
the state, the family, the courts and the schools as purely secular matters. That tension, Shelley
says, was the evangelical predicament throughout the Age of Ideologies.

\ssection{Return of the Old-Time Religion}
Carter, a former governor of Georgia, stood for a revived evangelical Christianity that took much of
the country by surprise. Revivalism had been pictured as the world of tent healers and frauds, loud,
tasteless and naive. Now politicians, athletes and entertainers talked openly about conversion, and
the religion of the ``Bible Belt'' was visible in every region. A generation of growth had brought
the number of Evangelicals to some forty-five million. The chapter sorts American Protestants into
two camps. The mainline bodies of the National Council of Churches (Episcopal, Methodist,
Presbyterian, American Baptist, United Church of Christ), which Martin Marty calls public
Protestants, had thirty-six million members, of whom perhaps a third were evangelical in conviction.
Another thirty-three million Evangelicals stood outside the Council. Taken together these ``private
Protestants'' came close to the forty-six million Roman Catholics.

They were not of one mind about the nation's troubles. The fundamentalists among them were militant
conservatives who rejected every compromise with modern culture. The Pentecostals, or charismatics,
spoke of a ``baptism in the Holy Spirit'' and practiced tongues and healing; they included
Episcopalians and close to a million Catholics as well as traveling healers. The majority were
ordinary Protestants who trusted the Bible, held orthodox doctrine and expected each believer to make
a deliberate commitment to Christ, whether suddenly or over time, which is what ``born again'' means.

\begin{sidebox}{Who Are the Charismatics?}
The label has several uses. Some denominations are charismatic as a whole (the Assemblies of God are
an example), and many others, such as the Episcopal Church, contain a sizable charismatic minority. A
workable rule of thumb is this: Christians are charismatic when the more spectacular gifts of the
Spirit, among them tongues, healing and prophecy, belong regularly to their public worship and to
their private prayer.
\end{sidebox}

Numbers and publicity did not give the movement influence over the nation's thinking, and the reason
lies some eighty years back. The chapter's title comes from George M.~Marsden of Notre Dame, who
likens Evangelicals to immigrants without roots in a strange country, with the difference that they
had not moved: the country had changed around them. After the Civil War the Southern churches lived
with defeat, and many Northern revivalists adopted premillennialism, the expectation that Christ will
return soon and that the world will get worse until he does. D.~L.~Moody put the mood in a famous
image: ``I look on the world as a wrecked vessel. God has given me a lifeboat and said, `Save all you
can.'\,'' Conservatives left social reform to others, a retreat the historian Timothy Smith named the
Great Reversal, and in the 1880s and 1890s they gave their energy to Bible study and personal
holiness. Fundamentalism grew from this soil.

Summer Bible conferences on the second coming began with a small gathering at Swampscott,
Massachusetts, in July 1876. The Boston Baptist A.~J.~Gordon helped to arrange big prophecy
conferences at New York in 1878 and Chicago in 1886, and there the premillennial scheme took firm
shape: sin increases until the Antichrist is let loose for a final devastation, and then Christ
comes with his saints to reign on earth for a thousand years. The same meetings refined a special
version of it, dispensationalism. In its classic form everything turns on the last three
dispensations, or ages, of history. Jesus came offering to rule Israel as an actual king in
Jerusalem. His people refused, so that age broke off, or was put on hold. He then turned to the
Gentiles with a covenant of grace in place of the covenant of law, and the age of the church began.
Like every earlier age it is expected to end in failure. When Christ returns he will govern the
world for a millennium through Israel. The book lays the scheme out in a table called ``Three Crucial Ages of Dispensationalism,'' which is given here in shorter form.
\begin{sumtable}[Three crucial ages of dispensationalism]{L{0.17\linewidth}L{0.22\linewidth}L{0.27\linewidth}L{0.22\linewidth}}
 & \sumhead{Age of Israel} & \sumhead{Age of the church} & \sumhead{Millennium} \\ \midrule
God's people & Israel & The church & Israel, with the nations \\
Made up mainly of & Hebrews & Gentiles & Hebrews, with the nations \\
What Jesus does & Offers himself as Israel's king & Refused by Israel, he goes to the Gentiles and
dies for sin & Comes again and reigns a thousand years \\
\end{sumtable}
Shelley regards the scheme as a sharp break with older premillennialism. It needs two peoples of
God, Israel and the church, living under two separate covenants and two principles, law and grace. In
its first form it therefore gave up several things Christians had long believed: that the promises of
the Old Testament are fulfilled in Jesus and his church, that Scripture is one book, that the Old
Testament knows grace and that the New demands obedience. It even handed large parts of the New
Testament, including the Sermon on the Mount, to Israel instead of the church. Later revisions made
the theory acceptable to more traditional believers, Presbyterians for instance, and it is now common
among conservative Protestants. The chapter notes an irony. Dispensationalists agree with liberal
scholars that Jesus first presented himself in earthly and political terms, and, like the doctrine
that miraculous gifts have ceased, their system accounts for the distance between church life today
and church life in the Bible.

Other Evangelicals of the period went to holiness conferences. Methodist revivalism had always taught
a ``second blessing,'' also called entire sanctification or Christian perfection, and the theme now
spread. New holiness bodies, the Church of the Nazarene among them, and ``deeper life'' meetings
called believers to surrender to the Spirit and so to live victoriously. Both interests, prophecy and
holiness, were biblical. Shelley adds that they also helped conservative Protestants to keep their
faith while losing their hold on the culture: people who could no longer direct public life found
consolation in the life of the Spirit.

\ssection{The Appearance of Fundamentalism}
The name fundamentalism goes back to \emph{The Fundamentals}, twelve paperback volumes of essays
defending basic doctrines, which appeared between 1910 and 1915. The idea and the money came from
Lyman Stewart, a California oil millionaire who believed that higher criticism and liberal theology
had to be answered, and from his brother Milton. Stewart heard Amzi C.~Dixon preach in August 1909
and engaged him for the project. Dixon, then pastor of Moody Church in Chicago, formed an editorial
committee on which the evangelist R.~A.~Torrey sat. The sixty-four contributors included many
premillennialists from America and Keswick teachers from England, and also conservative scholars of
other kinds, such as the Southern Baptist E.~Y.~Mullins and Princeton's Benjamin B.~Warfield. Three
million copies were mailed free of charge to pastors, missionaries and students of theology across
the world.

The open conflict with ``modernism'' waited until after the First World War and then broke out among
Baptists, Presbyterians, Methodists and Disciples of Christ. In 1920 Curtis Lee Laws, who edited the
Baptist \emph{Watchman-Examiner}, summoned the ``fundamentalists'' of the Northern Baptist Convention
to meet at Buffalo, and the moderate conservatives who formed the Fundamentalist Fellowship there
were the first to use the word of themselves. In their judgment the modernists were giving away the
core of the gospel: human sinfulness, the impossibility of salvation without grace, the death of
Christ as the source of new life for persons and for society, and the Bible as authoritative
revelation. They could not persuade the Convention to adopt a confession of faith. Laws nonetheless
wrote in 1924 that liberalism had been halted in some Baptist schools and that an inquiry into the
mission boards had achieved enough to make a separate mission needless. Harder-line Baptists thought
otherwise and organized the General Association of Regular Baptists.

Among Presbyterians the leading conservative was J.~Gresham Machen, a professor at Princeton Theological Seminary. The General Assembly approved in 1929 a reorganization of the seminary that combined its governing
boards. Machen and several eminent colleagues believed the change favored the liberals, resigned and
opened Westminster Theological Seminary in Philadelphia. Later he would not leave the Independent
Board for Presbyterian Foreign Missions, and a church court convicted him of defying his superiors.
The Orthodox Presbyterian Church and the Bible Presbyterian Church were the conservatives' reply.

The real dispute concerned the Bible. Shailer Mathews, dean of the University of Chicago Divinity
School, set out the liberal position in \emph{The Faith of Modernism} (1924) and claimed the support
of ``scientific investigations.'' Where the confessions treat Scripture as given from above, the
modernist reads it as ``the trustworthy record and product of a developing religion.'' To be a
Christian, on this view, is to share the attitudes and trust of the biblical writers, ``those ancient
men of imperfect morals whose hearts found God,'' and not to accept a body of writings. Shelley finds
the three marks of liberalism here: evolution applied to religion, a hopeful estimate of human nature
centered on religious experience, and a moral God who is easy to ``find.'' Machen, whom the book
considers the ablest spokesman of orthodoxy, had stated the other side in the \emph{Princeton
Theological Review} in 1915. The New Testament scholar, he wrote, must be a historian before anything
else, because the Bible reports an event, the life, death and resurrection of Jesus, which puts ``a
new face upon life.'' Its authority should be judged by whether it is right about him: was he a
teacher to imitate or a Savior? From this the book derives the marks of early fundamentalism: a
supernatural Christ vouched for by the resurrection, a reliable Bible as the source of faith, and a
humanity that needs to be made new.

\ssection{William Jennings Bryan, Symbol of Change}
These church quarrels reflected a change in the whole country, and the chapter traces it through the
career of William Jennings Bryan. He is remembered, through the play and film \emph{Inherit the
Wind}, for a few July days in 1925 at Dayton, Tennessee. John Scopes was on trial there for teaching
evolution in a state school, contrary to law, and Bryan clashed with the defense attorney Clarence
Darrow over the creation story in Genesis. Darrow was at his best, Bryan was old, and the reporters,
nearly all against him, wired millions of words about his discomfiture. The image of Bryan sweating
under cross-examination became a byword for the supposed stupidity of people who believe the Bible.
Shelley calls its reputation inflated and mostly undeserved and the film misleading: it casts Bryan
as a reactionary facing a man of progress, and it invents an episode in which he browbeats the
fianc\'ee of the accused.

The historical Bryan was the progressive. He spoke for ordinary people throughout his life, led the
reforming wing of the Democratic party, was three times its candidate for president and served
Woodrow Wilson as secretary of state. He fought social Darwinism and campaigned for woman suffrage,
two causes with a connection, since opponents of the vote claimed that women were too emotional, too
little evolved, to use it. As a man of peace he concluded arbitration treaties with thirty nations,
and he resigned in 1915 when Wilson and the cabinet refused to apply them after a German submarine
sank the \emph{Lusitania} on May~7. He then worked for Prohibition, which became law nationwide in
January 1920 through the Eighteenth Amendment. The book calls it probably the last evangelical
crusade for a moral America that succeeded. That later generations saw only sour puritanism in it
shows, Shelley thinks, how far the culture had moved. For Bryan it was one more battle against
interests that valued profit above people: ``It is not injustice to any man to refuse him permission
to enrich himself by injuring his fellowmen.''

The campaign against the teaching of evolution, his last, made him a fundamentalist leader. He had
worried about Darwinism for years; on reading \emph{The Descent of Man} in 1905 he predicted that its
view of human origins would undermine democracy. A set of speeches in the spring of 1921 placed him
at the head of the cause. His case was moral, not scientific. ``The objection to evolution is not,
primarily, that it is not true. The principal objection is that it is highly harmful to those who
accept it.'' His lecture ``The Menace of Darwinism,'' which circulated widely, reasoned that virtue
rests on belief in God, so that anything that erodes the belief erodes character. Applied to society,
evolution reduces human beings to animals, leaves no room for spiritual values and removes the
strongest motive for a good life; the strong may then prosper without a thought for the weak. In
1921--1929 twenty state legislatures received thirty-seven bills against evolution. Their sponsors
belonged to the generation that had passed Prohibition and trusted law to support morals.

Tennessee's statute forbade teachers in public schools to present any theory that contradicts the
biblical account of man's creation by God and to teach in its place descent from lower animals.
Whether Scopes, a young biology teacher, had done so hardly mattered once local promoters and the
press had turned the trial into a show. Bryan joined the prosecution, announcing that he had come
``to protect the Word of God against the greatest atheist or agnostic in the United States,'' to the
applause of the courtroom. Darrow replied that freedom of thought was the real issue, put Bryan in
effect on trial and used his answers to make fundamentalism look foolish. Scopes was convicted and
fined a nominal amount. Bryan carried Dayton, and Darrow carried the nation. Bryan died in his sleep
five days afterward, and in the book's judgment the evangelical campaign for a Christian America died
then too. The Depression followed. Conservative Protestantism lived on in its voluntary agencies and
on the radio, while public opinion fixed on a caricature of the fundamentalist as ignorant, narrow,
quarrelsome and separatist.

\ssection{The Renewal of Evangelical Christianity}
Evangelical religion came back into public view after the Second World War, and the voice most
people associated with it spoke with a soft North Carolina accent. Billy Graham grew up in Southern
Baptist and Southern Presbyterian surroundings, studied at three conservative schools (Bob Jones University, the Florida Bible Institute and finally Wheaton College), was briefly a pastor at Western Springs, Illinois, and then became an
evangelist for Youth for Christ. His Los Angeles meetings of 1949 made national news. Crusades in the
great stadiums of America and in cities overseas, a weekly radio program, network television and the
magazine \emph{Decision} made him known in almost every household.

Graham was the public face of a wider change. A number of conservatives had grown unhappy with the
direction fundamentalism took between the wars. They kept its doctrines and regretted its attitudes.
Carl F.~H.~Henry said so as early as 1947 in \emph{The Uneasy Conscience of Modern Fundamentalism},
which complained that the movement no longer brought Christian truth to bear on the great questions
of the age. In 1956 a group led by Graham founded the magazine \emph{Christianity Today}, and Henry
left Fuller Theological Seminary to edit it. Free copies went at first to ministers and seminarians
of every theological color, which gave it an immediate readership. According to Henry's first
editorial it meant to present historic Christianity to a new generation. Its founders held that
liberal theology had not answered people's moral and spiritual needs; its editors affirmed the
``complete reliability of the written Word of God'' and proposed to apply the gospel to the whole of
life.

Help came in 1960 from a direction nobody expected, Pentecostalism. The experience itself, ``the
baptism of the Holy Spirit with the evidence of speaking in tongues,'' was older. Isolated cases
aside, the worldwide movement started in 1906 with three years of revival at a Los Angeles mission on Azusa Street, to which visitors came from several continents and from which they took the message
home. Denominations formed, in the United States chiefly the Church of God in Christ, the Assemblies of God, the Pentecostal Holiness Church and the Church of God, and their members were mostly poor.
It was therefore news when, in 1960, tongues appeared in comfortable Lutheran and Episcopal parishes.
The newspapers spoke of neo-Pentecostalism. Charismatic prayer groups soon met all over the country,
made up largely of mainline Protestants who said their Christian life had gained a new dimension.

The chapter follows the renewal through two church families, as examples of what also happened in
other mainline churches and in Roman Catholicism. Among Lutherans it first stirred the American
Lutheran Church, where a dispute in 1963 calmed down once the leadership published guidelines. It
entered the Missouri Synod in 1967, and several pastors there lost their ministry over what they said
in public. In 1972 an international Lutheran conference on the Holy Spirit drew over ten thousand
people to Minneapolis. The United Presbyterian Church studied the question for two years, and the
report its General Assembly received in 1970, ``The Work of the Holy Spirit,'' was praised as one of
the best official treatments and encouraged further growth. By 1975 an estimated ten to fifteen
thousand members of the two large Presbyterian bodies had received the Pentecostal baptism.

The chapter ends in doubt. Evangelicals had long supposed that a large enough body of converted
people would change the moral tone of the country. At the end of the 1970s they had the numbers, and
crime, divorce, racism, violence, sexual disorder, alcoholism and drug abuse were as prevalent as
ever. Some of them asked whether Schlesinger's taunt applied to the whole movement, and whether
their faith, however deeply felt, made any public difference.

\begin{keyfacts}
\keypeople{Jimmy Carter; Dwight L.~Moody; A.~J.~Gordon; Lyman Stewart; Amzi C.~Dixon; Curtis Lee Laws;
J.~Gresham Machen; Shailer Mathews; William Jennings Bryan; Clarence Darrow; John Scopes; Billy
Graham; Carl F.~H.~Henry; George M.~Marsden.}
\keydates{%
  \kd{1878}{Prophecy conference in New York (Chicago followed in 1886)}
  \kd{1906}{Azusa Street revival begins in Los Angeles}
  \kd{1910--1915}{\emph{The Fundamentals} published}
  \kd{1920}{Baptist conservatives first call themselves ``fundamentalists''; Prohibition takes
    effect}
  \kd{1925}{Scopes trial at Dayton, Tennessee; death of Bryan}
  \kd{1929}{Machen leaves Princeton and founds Westminster Seminary}
  \kd{1949}{Billy Graham's Los Angeles campaign}
  \kd{1956}{\emph{Christianity Today} founded}
  \kd{1960}{Neo-Pentecostalism appears in mainline churches}
  \kd{1976}{Jimmy Carter's evangelical faith an issue in the presidential campaign}}
\keyterms{Evangelicals; born-again experience; mainline (public) and private Protestants;
premillennialism; dispensationalism; the Great Reversal; holiness movement; fundamentalism;
modernism; Prohibition; Pentecostalism; charismatics; neo-Pentecostalism.}
\end{keyfacts}
